\documentclass[11pt]{article}
\usepackage[margin=1in]{geometry}
\usepackage{amsmath,amssymb,amsthm}
\usepackage{booktabs}
\usepackage{enumitem}
\usepackage{algorithm}
\usepackage{algpseudocode}
\usepackage{hyperref}

\title{GQE: Reward Function and Training Pipeline}
\author{gqe-torch}
\date{}

\newcommand{\lammstruct}{\lambda_{\mathrm{struct}}}
\newcommand{\gammadepth}{\gamma_{D}}
\newcommand{\Lmax}{L_{\max}}
\newcommand{\Fhat}{\widehat{F}}

\begin{document}
\maketitle

\begin{abstract}
We describe the reward function and training pipeline of the JAX-based
Generative Quantum Eigensolver (GQE) in this repository. The reward is a
two-hyperparameter linear trade-off between process fidelity and a circuit
structure cost (CNOT count plus weighted depth), normalised by the maximum
allowed circuit length. Training is a GRPO-style policy-gradient loop over a
transformer that autoregressively emits variable-length circuits terminated by
a learned \textsc{stop} token.
\end{abstract}

\section{Reward function}

\subsection{Formulation}

For a sampled circuit $x$ with fidelity $F(x)\in[0,1]$, CNOT count $C(x)\in
\mathbb{Z}_{\ge 0}$, and depth $D(x)\in\mathbb{Z}_{\ge 0}$, the reward is
\begin{equation}
  R(x) \;=\; F(x) \;-\; \lammstruct\,\frac{C(x) + \gammadepth\, D(x)}{\Lmax},
  \label{eq:reward}
\end{equation}
where $\Lmax$ is the maximum gate budget (\texttt{model.max\_gates\_count},
a fixed architectural constant, not a tunable hyperparameter). The cost fed
into the GRPO objective is $-R(x)$.

\subsection{Hyperparameters}

The reward has exactly two hyperparameters, both exposed under the
\texttt{reward} section of the YAML config:

\begin{table}[h]
\centering
\begin{tabular}{lllp{7.3cm}}
\toprule
Symbol & Config key & Default & Role \\
\midrule
$\lammstruct$ & \texttt{reward.lambda\_structure} & $0.3$ &
Fidelity--structure trade-off. When $F$ is far from $1$, the fidelity term
dominates regardless of $\lammstruct$; as $F\to 1$ the structure penalty
becomes the discriminating signal between otherwise-equal circuits. \\
$\gammadepth$ & \texttt{reward.gamma\_depth} & $1.0$ &
Depth weight relative to CNOT count. Set to $0$ to optimise CNOT only;
raise above $1$ if depth matters more than CNOT for the target platform. \\
\bottomrule
\end{tabular}
\caption{Reward hyperparameters.}
\label{tab:reward-hp}
\end{table}

Additional reward-section config keys are \emph{not} hyperparameters of the
reward itself; they only configure the Pareto archive used for reporting:

\begin{table}[h]
\centering
\begin{tabular}{llp{8cm}}
\toprule
Key & Default & Use \\
\midrule
\texttt{reward.enabled} & \texttt{true} & Toggle the full reward + Pareto
archive. When false the cost signal reduces to $1-F$ (fidelity-only). \\
\texttt{reward.fidelity\_floor} & $0.0$ & Minimum $F$ for a circuit to
enter the Pareto archive. \\
\texttt{reward.max\_archive\_size} & $500$ & Crowding-distance pruning
threshold for the Pareto archive. \\
\texttt{reward.fidelity\_threshold} & $0.99$ & Threshold used by the
reporting queries \texttt{best\_by\_cnot}/\texttt{best\_by\_depth}
(``best structure at $F\ge$ threshold''). \\
\bottomrule
\end{tabular}
\caption{Reporting / archive configuration (not in the reward formula).}
\end{table}

\subsection{Behavioural properties}

The linear reward (\ref{eq:reward}) is chosen because, despite its simplicity,
it aligns the policy gradient with all three objectives the user cares about:
maximum fidelity, minimum CNOT count, and minimum depth.

\begin{description}
  \item[Fidelity dominates when it is poor.] Since $F\in[0,1]$ and
  $C+\gammadepth D \le (1+\gammadepth)\Lmax$, the structure term is bounded
  by $\lammstruct(1+\gammadepth)$, which is typically $\ll 1$. A circuit
  with fidelity $0.9$ beats a circuit with fidelity $0.5$ in reward units
  regardless of structure.
  \item[Structure discriminates only near $F=1$.] When two circuits have
  comparable fidelity, the term $-\lammstruct(C+\gammadepth D)/\Lmax$ breaks
  the tie in favour of the shorter, CNOT-lighter one. This reproduces the
  lexicographic preference ``fidelity first, then CNOT, then depth'' as a
  smooth gradient signal with no regime switches.
  \item[Trivial circuits are not rewarded.] The minimum-length circuit
  (one gate) has $C+\gammadepth D \le 1+\gammadepth$, so the structure
  bonus is tiny; low-fidelity short circuits lose to high-fidelity long ones.
  A $0$-gate circuit is disallowed by construction (STOP is masked at the
  first sampled position).
  \item[Scale-invariant in $\Lmax$.] Dividing by $\Lmax$ makes the structure
  term have the same magnitude across different maximum-length settings,
  so $\lammstruct$ does not need to be retuned when $\Lmax$ changes.
\end{description}

\subsection{What was removed}

The previous reward (documented in \texttt{gqe\_reward\_design.tex}) was a
piecewise function of the log-infidelity gap $\Delta_\Phi$ relative to a
reference circuit, gated by $w(F)=\sigma((F-F_{\mathrm{gate}})/\tau)$ and
split into ``bad drop / mild drop / improvement'' regimes. It exposed
approximately twenty hyperparameters
($\varepsilon_\phi$, $F_{\mathrm{gate}}$, $\tau$, $\delta_{\mathrm{bad}}$,
$\alpha_{\mathrm{bad/soft/good}}$, $\kappa_{\mathrm{bad/soft/good}}$,
$\lambda_D$, $\lambda_C$, $\mu_D$, $\mu_C$, $\eta$, $\alpha_{\mathrm{abs}}$,
$\beta_1$, $\beta_2$, the reference EMA rate, the reference mode, the late
fidelity floor, its ramp epoch, and the warmup rescore length). All of
these are gone in the current design, along with the running reference
tracker and the Pareto-proxy bonus.

\section{Training pipeline}

\subsection{Overview}

GQE trains a GPT-2-style transformer $\pi_\theta$ to autoregressively emit
circuit token sequences that compile to a target unitary $U_{\mathrm{target}}$.
The pipeline is a GRPO policy-gradient loop:

\begin{enumerate}[leftmargin=1.8em]
  \item Sample a batch of circuits from $\pi_\theta$ using KV-cached
  autoregressive decoding (Sec.~\ref{sec:rollout}).
  \item Optionally refine each circuit's rotation angles with L-BFGS
  (Sec.~\ref{sec:contopt}).
  \item Score each circuit with the reward (\ref{eq:reward}); push
  $(\text{tokens}, \text{cost}, \log\pi_{\theta_{\mathrm{old}}})$ into a
  replay buffer.
  \item Take gradient steps on the clipped GRPO objective using
  length-masked sequence log-probabilities (Sec.~\ref{sec:grpo}).
  \item Repeat.
\end{enumerate}

\subsection{Gate pool and token vocabulary}

The gate pool is built by \texttt{operator\_pool.build\_operator\_pool}:
for each qubit, one token per configured rotation axis
(\texttt{pool.rotation\_gates}, any subset of \{RX, RY, RZ\}) plus one
\textsc{sx} token; for each ordered pair $(c,t)$, one \textsc{cnot} token.
Two special tokens are prepended to form the transformer vocabulary:
\[
\mathcal{V} \;=\; \{\underbrace{\text{\textsc{bos}}}_{\text{id}=0},
                   \underbrace{\text{\textsc{stop}}}_{\text{id}=1},
                   g_1, g_2, \ldots, g_{|\mathrm{pool}|}\},
\qquad |\mathcal{V}| = |\mathrm{pool}| + 2.
\]
\textsc{bos} always occupies position $0$ of every sequence; \textsc{stop}
terminates a variable-length circuit.

\subsection{Rollout and \textsc{stop}-token sampling}
\label{sec:rollout}

At each autoregressive step the model emits logits $\ell_t\in\mathbb{R}^{|
\mathcal{V}|}$ and samples the next token from $\mathrm{softmax}(-\beta\ell_t)$,
where $\beta$ is the inverse temperature from the scheduler.

Two position-dependent action masks are applied before sampling:
\begin{itemize}[leftmargin=1.5em]
  \item \textsc{bos} is unsamplable at every step.
  \item \textsc{stop} is unsamplable at step $t=0$ only, so every circuit has
  at least one gate.
\end{itemize}
Once a sample emits \textsc{stop}, every subsequent emitted token is
force-replaced with \textsc{stop}. The resulting sequence has shape
$(1+\Lmax)$: a \textsc{bos} prefix, a prefix of real gate tokens, a single
\textsc{stop} token (if terminated), and a tail of \textsc{stop} padding.

The KV-cache rollout in \texttt{kv\_rollout.build\_kv\_rollout\_fn} runs in
$O(\Lmax)$ transformer work per circuit and is the hot-path generator used
during training. A fallback full-attention sampler in
\texttt{pipeline.\_generate\_rollout} is used for arbitrary-prefix generation.

\paragraph{Length derivation.} Given a sampled token sequence
$x_{1:\Lmax}$ (\textsc{bos} stripped), the per-sample length
\[
  L(x) \;=\; \min\!\bigl(\,\arg\min_t\!\{t \mid x_t = \textsc{stop}\} + 1,\;
  \Lmax\bigr)
\]
is computed by \texttt{kv\_rollout.compute\_lengths\_from\_tokens}. $L(x)$
counts \textsc{stop} itself as a real decision (the model gets credit for
deciding to stop) but excludes all forced-\textsc{stop} padding.

\subsection{Circuit metrics with \textsc{stop} padding}

Fidelity, depth, CNOT count, and total gate count are all computed on the
full fixed-length sequence, treating \textsc{bos} and \textsc{stop} tokens
as no-ops:
\begin{itemize}[leftmargin=1.5em]
  \item The token-matrix table \texttt{pipeline.token\_matrices\_jax} is
  vocab-indexed; the \textsc{bos} and \textsc{stop} entries are the identity
  matrix, so $U_{\mathrm{circuit}}(x) = \prod_{t=1}^{\Lmax} M_{x_t}$
  collapses the padding to identities.
  \item The depth/total-gates scan (\texttt{pipeline.\_make\_batch\_structure\_fn})
  consults a per-token \texttt{token\_is\_noop} mask so \textsc{stop} tokens
  do not advance per-qubit depth.
  \item The CNOT count uses \texttt{two\_qubit\_token\_mask} which is
  \texttt{False} for both \textsc{bos} and \textsc{stop}.
\end{itemize}

\subsection{Continuous angle optimisation}
\label{sec:contopt}

When \texttt{continuous\_opt.enabled=true}, the L-BFGS angle optimiser in
\texttt{continuous\_optimizer.ContinuousOptimizer} takes the sampled gate
\emph{structure} (token ids) and optimises the rotation angles to maximise
fidelity. The vocabulary-sized encoding tables map \textsc{bos} and
\textsc{stop} to a \textsc{noop} gate type that contributes identity to the
circuit unitary, so the optimiser processes the full fixed-length batch
without ragged arrays.

Key config keys:
\begin{itemize}[leftmargin=1.5em]
  \item \texttt{continuous\_opt.optimizer}: \texttt{lbfgs} (default) or
  \texttt{adam}.
  \item \texttt{continuous\_opt.steps}: maximum inner iterations per circuit.
  \item \texttt{continuous\_opt.num\_restarts}: independent random restarts
  (best-of-$k$).
  \item \texttt{continuous\_opt.top\_k}: if $>0$, only the top-$k$ circuits
  by discrete fidelity within each rollout are refined; the rest keep their
  discrete fidelity.
\end{itemize}

\subsection{GRPO objective with length masking}
\label{sec:grpo}

Let $\theta_{\mathrm{old}}$ be the policy at sampling time and $\theta$ be the
current training-time policy. For a batch of sampled circuits with costs
$c_i = -R(x_i)$, GRPO optimises
\begin{equation}
  \mathcal{L}(\theta) \;=\; -\frac{1}{N}\sum_i \min\!\Bigl(
    r_i(\theta)\,A_i,\;
    \mathrm{clip}\!\bigl(r_i(\theta),1\!-\!\varepsilon,1\!+\!\varepsilon\bigr)\,A_i
  \Bigr),
\end{equation}
with advantage $A_i = (\bar{c} - c_i)/\mathrm{std}(c)$ and importance ratio
$r_i(\theta) = \exp\!\bigl(\bar{\ell}_i(\theta) - \bar{\ell}_i(\theta_{
\mathrm{old}})\bigr)$. The clip width $\varepsilon$ is
\texttt{training.grpo\_clip\_ratio}.

The sequence log-probability is a \emph{length-masked mean} over valid
(non-padding) positions:
\begin{equation}
  \bar{\ell}_i(\theta) \;=\; \frac{1}{L(x_i)} \sum_{t=1}^{L(x_i)}
  \log \pi_\theta(x_i^{(t)} \mid x_i^{(0:t-1)}),
  \label{eq:length-mean}
\end{equation}
implemented in \texttt{loss.reduce\_sequence\_log\_probs}. The
length-normalised mean (rather than a sum) keeps importance ratios
comparable across circuits of different lengths, and the mask
(\texttt{positions < length}) prevents forced-\textsc{stop} padding from
contaminating the gradient. The same action masks used at sampling
(\textsc{bos} always, \textsc{stop} at position~0) are re-applied via
\texttt{loss.apply\_action\_masks} when recomputing log-probabilities, so
the behaviour and training policies share a consistent support.

\subsection{Replay buffer and training schedule}

The replay buffer (\texttt{data.ReplayBuffer}) is a FIFO deque of
$(\text{tokens}, \text{cost}, \log\pi_{\theta_{\mathrm{old}}})$ triples of
capacity \texttt{buffer.max\_size}. Each epoch:
\begin{enumerate}[leftmargin=1.8em]
  \item \texttt{collect\_rollout}: sample \texttt{training.num\_samples}
  circuits, score them, push to the buffer, update the Pareto archive,
  update the temperature scheduler with fidelity costs $1-F$.
  \item Train for \texttt{buffer.steps\_per\_epoch} passes over the buffer
  using \texttt{training.batch\_size}-sized shuffled minibatches.
\end{enumerate}
There is no separate warmup phase: because the reward (\ref{eq:reward}) is
smooth and already fidelity-dominated when $F$ is low, training uses the
full reward from epoch $0$.

\subsection{Temperature scheduler}

The inverse temperature $\beta$ is controlled by
\texttt{factory.create\_temperature\_scheduler} and can be one of
\texttt{fixed}, \texttt{linear}, or \texttt{cosine} (config
\texttt{temperature.scheduler}). The scheduler's \texttt{update} receives the
raw fidelity costs $1-F$ (not the scalarised reward) so its annealing is not
confused by changes in reward scale.

\subsection{Pareto archive (reporting only)}

\texttt{pareto.ParetoArchive} maintains the non-dominated front over
$(F\uparrow, D\downarrow, C\downarrow)$ with crowding-distance pruning. It is
updated every rollout with every sampled circuit. In the current design it
does not feed back into the reward; it supplies reporting queries
(\texttt{best\_by\_fidelity}, \texttt{best\_by\_cnot}, \texttt{best\_by\_depth})
at end of training.

\subsection{Algorithm summary}

\begin{algorithm}[H]
\caption{GQE training epoch}
\begin{algorithmic}[1]
\State \textbf{Inputs:} policy $\pi_\theta$, target $U_{\mathrm{target}}$,
replay buffer $\mathcal{B}$, scheduler $\mathrm{sch}$,
Pareto archive $\mathcal{A}$.
\State $\beta \leftarrow \mathrm{sch.get}()$
\State $\{x_i\}_{i=1}^{N} \leftarrow \textsc{KVRollout}(\pi_\theta, \beta)$
\Comment{\textsc{bos}-prefixed, \textsc{stop}-padded, shape $N\!\times\!(1+\Lmax)$}
\State $\{L(x_i)\} \leftarrow \textsc{LengthsFromTokens}(\{x_i\})$
\ForAll{$x_i$}
  \State $F_i \leftarrow \textsc{Fidelity}(x_i;\,U_{\mathrm{target}})$
  \Comment{L-BFGS-refined if enabled}
  \State $D_i, C_i \leftarrow \textsc{DepthAndCNOT}(x_i)$
  \Comment{\textsc{stop} tokens skipped}
  \State $c_i \leftarrow -\bigl(F_i - \lammstruct(C_i + \gammadepth D_i)/\Lmax\bigr)$
  \State $\ell^{\mathrm{old}}_i \leftarrow
  \bar{\ell}_i(\theta) = \tfrac{1}{L(x_i)}\sum_{t=1}^{L(x_i)}\log\pi_\theta(x_i^{(t)} \mid x_i^{(<t)})$
  \State $\mathcal{B}.\mathrm{push}(x_i, c_i, \ell^{\mathrm{old}}_i)$
  \State $\mathcal{A}.\mathrm{update}(\mathrm{ParetoPoint}(F_i, D_i, C_i, x_i))$
\EndFor
\State $\mathrm{sch.update}(\{1-F_i\})$
\For{$k = 1 \dots \mathrm{steps\_per\_epoch}$}
  \ForAll{minibatch $B \subset \mathcal{B}$}
    \State $\theta \leftarrow \theta - \eta \nabla_\theta
    \mathcal{L}_{\mathrm{GRPO}}(\theta; B, \beta)$
    \Comment{clipped surrogate, length-masked $\bar\ell$}
  \EndFor
\EndFor
\end{algorithmic}
\end{algorithm}

\subsection{Configuration hyperparameters at a glance}

\begin{table}[h]
\centering
\small
\begin{tabular}{llp{8cm}}
\toprule
Section & Key & Role \\
\midrule
\texttt{target} & \texttt{num\_qubits}, \texttt{type}, \texttt{path} &
Target unitary specification. \\
\texttt{pool} & \texttt{rotation\_gates} & Rotation axes included in the
gate pool (subset of \{rx, ry, rz\}). \\
\texttt{model} & \texttt{size} & \texttt{tiny|small|medium|large} — GPT-2
width/depth preset. \\
\texttt{model} & \texttt{max\_gates\_count} & $\Lmax$: maximum emitted
circuit length. \\
\texttt{training} & \texttt{max\_epochs}, \texttt{num\_samples},
\texttt{batch\_size}, \texttt{lr}, \texttt{grad\_norm\_clip} &
Standard optimisation controls. \\
\texttt{training} & \texttt{grpo\_clip\_ratio} & $\varepsilon$ in the
clipped surrogate. \\
\texttt{training} & \texttt{early\_stop} & Terminate on $F\ge 1 - 10^{-8}$. \\
\texttt{temperature} & \texttt{scheduler}, \texttt{initial\_value},
\texttt{delta}, \texttt{min\_value}, \texttt{max\_value} &
Inverse-temperature $\beta$ schedule for sampling. \\
\texttt{buffer} & \texttt{max\_size}, \texttt{steps\_per\_epoch} &
Replay FIFO capacity and number of training passes per rollout. \\
\texttt{continuous\_opt} & \texttt{enabled}, \texttt{optimizer},
\texttt{steps}, \texttt{lr}, \texttt{top\_k}, \texttt{num\_restarts} &
L-BFGS/Adam angle-refinement controls. \\
\texttt{reward} & \texttt{lambda\_structure}, \texttt{gamma\_depth} &
\textbf{The two reward hyperparameters.} See Table~\ref{tab:reward-hp}. \\
\texttt{reward} & \texttt{fidelity\_floor}, \texttt{max\_archive\_size},
\texttt{fidelity\_threshold} & Pareto archive (reporting). \\
\bottomrule
\end{tabular}
\caption{Full configuration map.}
\end{table}

\end{document}
